\loai{Tính chu kì - hằng số phóng xạ}
\begin{vidu}%[3P7-7.2Y][Loai1]
Hạt nhân $\Hn{90}{227}{Th}$ là phóng xạ $\alpha$ có chu kì bán rã là 18,3 ngày. Hằng số phóng xạ của hạt nhân là
\choice
{\True $4,38.10^{-7}\ s^{-1}$}
{$0,038\ s^{-1}$}
{$26,4\ s^{-1}$}
{$0,0016\ s^{-1}$}
\loigiai{
\begin{itemize}
\item Chu kì bán rã: $T=18,3.86400=1581120\ s$.
\item Hằng số phóng xạ: $\lambda = \dfrac{\ln 2}{T}=4,38.10^{-7}\ s^{-1}$. 
\end{itemize}
}
\end{vidu}
\loai{Xác định loại phóng xạ}
\begin{vidu}%[3P7-7.2Y][Loai2]
\nguon{ĐH - 2008} Hạt nhân $\Hn{88}{226}{Ra}$ biến đổi thành hạt nhân $\Hn{86}{222}{Rn}$ do phóng xạ
\choice
{$\alpha$ và $\beta^-$}
{$\beta^-$}
{\True $\alpha$}
{$\beta^+$}
\loigiai{
\begin{itemize}
\item Phương trình phóng xạ: 
$\Hn{88}{226}{Ra}\rightarrow \Hn{86}{222}{Ra}+^A_ZX \Rightarrow
\begin{cases}
A=226-222=4\\
Z=88-86=2
\end{cases}$
\item Vậy X là hạt $\alpha$. 
\end{itemize}
}
\end{vidu}
\loai{Xác định hạt nhân con}
\begin{vidu}%[3P7-7.2B][Loai3]
Trong quá trình phân rã hạt nhân $\Hn{92}{238}{U}$ thành hạt nhân $\Hn{92}{234}{U}$, đã phóng ra một hạt $\alpha$ và hai hạt
\choice
{neutron}
{\True electron}
{positron}
{proton}
\loigiai{
\begin{itemize}
\item Phương trình phóng xạ: 
$\Hn{92}{238}{U}\rightarrow \Hn{92}{234}{Ra}+_2^4\alpha+2^A_ZX \Rightarrow
\begin{cases}
A=\dfrac{238-234-4}{2}=0\\
Z=\dfrac{92-92-2}{2}=-1
\end{cases}$
\item Vậy X là hạt $\beta^-$ (electron). 
\end{itemize}
}
\end{vidu}
\begin{vidu}%[3P7-7.2B][Loai3]
$\Hn{83}{210}{Bi}$ (bismut) là chất phóng xạ $\beta^-$. Hạt nhân con (sản phẩm của phóng xạ) có cấu tạo gồm
\choice
{84 neutron và 126 prôton}
{\True 126 neutron và 84 prôton}
{83 neutron và 127 prôton}
{127 neutron và 83 prôton}
\loigiai{
\begin{itemize}
\item Phương trình phóng xạ: 
$\Hn{83}{210}{U}\rightarrow _{-1}^0\beta^-+^A_ZX \Rightarrow
\begin{cases}
A=210-0=210\\
Z=83+1=84
\end{cases}$
\item Vậy X có $\begin{cases}
84\ \ct{proton} \\
210-84=126\ \ct{neutron}
\end{cases} $. 
\end{itemize}
}
\end{vidu}
\begin{vidu}%[3P7-7.2B][Loai3]
Chất radi phóng xạ $\alpha$ có phương trình: $\Hn{88}{226}{Ra}\to \alpha + \Hn{x}{y}{Ra}$
\choice
{\True $x=86$; $y=222$}
{$x=84$; $y=222$}
{$x=224$; $y=84$}
{$x=224$; $y=86$}
\loigiai{
Áp dụng định luật bảo toàn số khối và điện tích ta được:
\begin{align*}
88=2+x\Rightarrow x = 86\,\,\,\text{và}\,\,\,226=4+y\Rightarrow y = 222.
\end{align*}
}
\end{vidu}
\begin{vidu}%[3P7-7.2Y][Loai3]
Hạt nhân $\Hn{6}{11}{Cd}$ phóng xạ $\beta^+$, hạt nhân con là
\choice
{$\Hn{7}{11}{N}$}
{\True $\Hn{5}{11}{B}$}
{$\Hn{8}{15}{O}$}
{$\Hn{7}{12}{N}$}
\loigiai{
Ta có: $\Hn{6}{11}{Cd}\to _{1}^{0}{{\beta}^{+}}+_{6-1}^{11}X\Rightarrow _{6-1}^{11}X=_{5}^{11}X\equiv \Hn{5}{11}{B}$.
}
\end{vidu}
\begin{vidu}%[3P7-7.2Y][Loai3]
Hạt nhân $\Hn{6}{14}{C}$ phóng xạ $\beta^-$. Hạt nhân con sinh ra có
\choice
{5p và 6n}
{6p và 7n}
{\True 7p và 7n}
{7p và 6n}
\loigiai{
Ta có: $\Hn{6}{14}{C}\to\Hn{{-1}}{0}{{\beta}^{-}}+\Hn{{6-(-1)}}{14}{X}\Rightarrow \Hn{{6-(-1)}}{14}{X}=\Hn{7}{7}{X}\Rightarrow \begin{cases}
7\ \ct{proton} \\
7\ \ct{neutron}
\end{cases}$
}
\end{vidu}
\begin{vidu}%[3P7-7.2B][Loai3]
Hạt nhân $\Hn{84}{210}{Po}$ phân rã $\alpha$ thành hạt nhân con X. Số nuclôn trong hạt nhân X bằng
\choice
{82}
{210}
{124}
{\True 206}	
\loigiai{
\begin{itemize}
\item Phương trình phản ứng: $\Hn{84}{210}{Po}\to_2^4\alpha +\Hn{82}{206}{X}.$
\item Vậy hạt nhân X có 206 nuclôn.
\end{itemize}
}
\end{vidu}
\loai{Xác định tên hạt nhân sau chuỗi phóng xạ}
\begin{vidu}%[3P7-7.2B][Loai4]
Hạt nhân $\Hn{88}{226}{Ra}$ phóng ra 3 hạt $\alpha$ và một hạt $\beta^-$ trong chuỗi phóng xạ liên tiếp. Khi đó hạt nhân con tạo thành là	
\choice
{$_{84}^{222}X$}
{\True $_{83}^{214}X$}
{$_{83}^{222}X$}
{$_{84}^{224}X$}	
\loigiai{
\begin{itemize}
\item Phương trình phóng xạ: 
$\Hn{88}{226}{Ra}\rightarrow 3^4_2He+_{-1}^0\beta^-+^A_ZX \Rightarrow
\begin{cases}
A=214\\
Z=83
\end{cases}$
\end{itemize}
}
\end{vidu}
\begin{vidu}%[3P7-7.2B][Loai4]
Hạt nhân $\Hn{91}{234}{Pa}$ phóng xạ $\beta^-$ tạo thành hạt nhân X. Hạt nhân X tiếp tục phóng xạ $\alpha$ tạo thành
\choice
{$\Hn{92}{234}{U}$}
{$\Hn{88}{230}{Ra}$}
{$\Hn{90}{234}{U}$}
{\True $\Hn{90}{230}{Th}$}
\loigiai{
\begin{itemize}
\item Hạt nhân $\Hn{91}{234}{Pa}$ phóng xạ $\beta^-$ tạo thành hạt nhân X
\begin{align*}
\Hn{91}{234}{Pa}\to \Hn{{ - 1}}{0}{\beta}^ - + _{Z}^{A}X \Rightarrow 
\begin{cases}
A=234\\
Z=92
\end{cases} 
\end{align*}
\item Hạt nhân $\Hn{90}{234}{X}$ phóng xạ $\alpha$ tạo thành hạt nhân Y
\begin{align*}
\Hn{90}{234}{X}\to \Hn{2}{4}{\alpha} + _{Z'}^{A'}Y \Rightarrow 
\begin{cases}
A'=230\\
Z'=90
\end{cases} 
\end{align*}
\end{itemize}
}
\end{vidu}
\begin{vidu}%[3P7-7.2B][Loai4]
Hạt nhân $\Hn{92}{238}{U}$ sau 8 phân rã $\alpha$ và 6 phân rã $\beta^-$ sẽ biến đổi thành hạt nhân
\choice
{$\Hn{88}{226}{Ra}$}
{\True $\Hn{82}{206}{Pb}$}
{$\Hn{83}{210}{Bi}$}
{$\Hn{84}{210}{Po}$}	
\loigiai{
\begin{itemize}
\item Phản ứng hạt nhân: $\Hn{92}{238}{U}\to 8\alpha +6\beta^-+\Hn{Z}{A}{X}$.
\begin{itemize}
\item Bảo toàn số khối: $238=8.A_{He}+A\Rightarrow A=238-8.A_{He}=2388.4=206$.
\item Bảo toàn điện tích, ta có $92=8.Z_{He}+6Z_e+Z\Rightarrow Z=92-8Z_{He}-6Z_e=92-8.2-6.\left({-1}\right)=82$.
\end{itemize} 
\item Vậy X là $\Hn{82}{206}{Pb}$.
\end{itemize}	
}
\end{vidu}
\loai{Xác định số lần phóng xạ trong chuỗi phóng xạ}
\begin{vidu}%[3P7-7.2B][Loai5]
Đồng vị $\Hn{92}{234}{U}$ sau một chuỗi phóng xạ $\alpha$ và $\beta^-$ biến đổi thành $\Hn{82}{206}{Pb}$. Số phóng xạ $\alpha$ và $\beta^-$ trong chuỗi là
\choice
{\True 7 phóng xạ $\alpha$ , 4 phóng xạ $\beta^-$}
{5 phóng xạ $\alpha$ , 5 phóng xạ $\beta^-$}
{10 phóng xạ $\alpha$ , 8 phóng xạ $\beta^-$}
{16 phóng xạ $\alpha$ , 12 phóng xạ $\beta^-$}
\loigiai{
\begin{itemize}
\item Phương trình phóng xạ: 
\begin{align*}
\Hn{92}{234}{U}\rightarrow x_2^4He + y_{ - 1}^0{{\beta}^{ - }} + \Hn{82}{206}{Pb} \Rightarrow\begin{cases}
234=4x + 206\\
92=2x - y + 82
\end{cases} \Rightarrow\begin{cases}
x=7\\
y=4
\end{cases}
\end{align*}
\item Vậy số phóng xạ $\alpha$ và $\beta^-$ trong chuỗi là 7 phóng xạ $\alpha$ , 4 phóng xạ $\beta^-$.
\end{itemize}
}
\end{vidu}
\begin{vidu}%[3P7-7.2B][Loai5]
Sau bao nhiêu lần phóng xạ $\alpha$ và bao nhiêu lần phóng xạ $\beta$ cùng loại thì hạt nhân thori $\Hn{90}{232}{Th}$ biển đổi thành hạt nhân chì $\Hn{82}{208}{Pb}$?
\choice
{\True 6 lần phóng xạ $\alpha$ và 4 lần phóng xạ $\beta^-$}
{6 lần phóng xạ $\alpha$ và 4 lần phóng xạ $\beta^+$}
{8 lần phóng xạ $\alpha$ và 6 lần phóng xạ $\beta^-$}
{8 lần phóng xạ $\alpha$ và 6 lần phóng xạ $\beta^+$}
\loigiai{
$\Hn{90}{232}{Th}\rightarrow x\Hn{2}{4}{He}+y\Hn{{-1}}{0}{{\beta}}+\Hn{82}{208}{Pb}$
$\Rightarrow \begin{cases}232=4x+208 \\
90=2x-y+82 
\end{cases} $
$\Rightarrow \begin{cases} x=6 \\
y=4 
\end{cases} $ 
$\Rightarrow$ 6 phóng xạ $\alpha$ , 4 phóng xạ $\beta^-$.
}
\end{vidu}